% ---------------------------------------------------------------------
%  Introduction (non numérotée) : but + planification de l'expérience
%  (couvre les points « but du stage » et « planification » des consignes)
% ---------------------------------------------------------------------
\chapter*{Introduction}
\addcontentsline{toc}{chapter}{Introduction}
\label{chap:intro}

\section*{Une note sur le vocabulaire}
\label{sec:vocabulaire}

Le domaine financier emploie un vocabulaire largement anglophone, y compris
dans les travaux rédigés en français. Ces termes sont ici traduits de façon
systématique, le mot anglais consacré étant rappelé entre parenthèses lors de
sa première apparition. Font exception les noms propres de méthodes
statistiques, qui restent sous la forme sous laquelle la littérature les
désigne, ainsi qu'une variable centrale de ce travail, l'\edgevar{}, dont le
nom est conservé pour éviter toute ambiguïté avec le sens courant du mot
français correspondant.

\section*{But du mémoire}
\label{sec:but}

% ANGLE DU MÉMOIRE (décidé) : ce n'est PAS une étude de stratégies de
% trading, mais la construction d'un OUTIL DE VALIDATION statistique.
% Les stratégies ne sont que des jeux de données de test.

Lors de mon stage en entreprise, j'ai reçu pour consigne de créer un agent
d'intelligence artificielle capable d'intervenir en bourse. Un tel agent
achète et revend des titres, c'est-à-dire des parts de propriété de sociétés
cotées, en suivant des règles de décision qu'on lui fournit. J'ai choisi de
doter cet agent de capacités d'analyse statistique. Après avoir construit ce
qui lui permettait de passer des ordres d'achat et de vente, et d'accéder aux
données du marché, il m'a semblé qu'il lui manquait l'essentiel, à savoir la
capacité d'évaluer lui-même les règles de décision qu'il emploie.

Ce rapport présente l'outil d'évaluation que je lui ai donné.

L'objectif de cet outil n'est pas de découvrir une règle de décision
rentable. Il est de construire un enchaînement de méthodes capable de
déterminer, pour une stratégie quelconque qu'on lui soumet, si l'avantage
qu'elle affiche est réel, ou s'il relève du hasard de l'échantillon ou de la
tendance générale du marché. Une stratégie désigne ici une règle fixée à
l'avance, qui indique à quel moment acheter un titre et à quel moment le
revendre.

La contribution revendiquée est donc méthodologique. Elle porte sur la chaîne
de tests elle-même, sur sa justification et sur la vérification de ses
conditions d'application.

Les dix stratégies d'analyse technique employées dans ce rapport ne
constituent pas l'objet d'étude, mais un banc d'essai. L'analyse technique
regroupe les méthodes qui cherchent à anticiper l'évolution d'un cours à
partir de la seule lecture de son historique de prix, sans considération pour
la santé économique de l'entreprise. Ces stratégies fournissent des jeux de
données réels, qui permettent de montrer que l'outil discrimine correctement,
c'est-à-dire qu'il sait rejeter ce qui ne tient pas et retenir ce qui résiste
à l'analyse. Les résultats obtenus sont donc présentés à titre d'illustration
du bon fonctionnement du protocole, et non comme des recommandations
d'investissement.

Pour s'assurer que l'outil ne rejette pas tout par excès de sévérité, on lui
soumet également un témoin, une stratégie dont l'avantage est réel par
construction.

\section*{Les trois menaces qui pèsent sur le jugement}
\label{sec:menaces}

Trois difficultés se dressent devant quiconque veut juger une stratégie sur
des données de marché. Aucune n'est propre à ce travail, mais chacune suffit
à fausser un verdict, et c'est leur traitement conjoint qui fait l'intérêt du
protocole.

La première tient à la dérive du marché. Le premier réflexe pour juger une
stratégie consiste à examiner si elle rapporte, un gain positif semblant
suffire à la déclarer efficace. Ce raisonnement est pourtant trompeur, car
encore faut-il que ce gain dépasse ce qu'une entrée au hasard aurait
rapporté.

Les stratégies sont ici appliquées aux titres du S\&P~500, l'indice qui
regroupe les plus grandes capitalisations américaines, c'est-à-dire les
sociétés dont la valeur totale en bourse est la plus élevée. Or cet indice
progresse sur toute la période étudiée. Cette dérive haussière pose un
problème, puisqu'une stratégie peut réaliser des gains sans avoir capté le
moindre signal, simplement portée par la montée générale. C'est pourquoi on
mesurera non pas le gain brut, mais l'avantage propre qu'apporte chaque
stratégie par rapport au hasard.

La deuxième difficulté provient du nombre de tests. Plus on examine de
stratégies, plus le risque d'en voir au moins une paraître gagnante par pur
hasard augmente. Cette difficulté sera traitée au
chapitre~\ref{chap:juger} par la correction de Bonferroni.

La troisième, sans doute la plus délicate, est la dépendance des données. De
nombreuses méthodes statistiques exigent que les observations soient
indépendantes les unes des autres, chacune apportant une information neuve.
Or les opérations réalisées par une stratégie se ressemblent sous plusieurs
aspects, qu'il s'agisse de la période durant laquelle elles se déroulent, du
titre sur lequel elles portent ou du secteur d'activité de celui-ci. Cette
difficulté sera mesurée au chapitre~\ref{chap:donnees}, puis neutralisée pas
à pas au chapitre~\ref{chap:juger}.

Ignorer une seule de ces trois menaces conduit mécaniquement à un verdict
trop optimiste. Les traiter ensemble est précisément ce qui distingue un
outil de validation d'un simple calcul de rentabilité.

\section*{Démarche du mémoire}
\label{sec:planification}

Le mémoire suit le fil d'une enquête unique, du diagnostic au verdict, puis
du verdict à l'outil. Il commence par le terrain et ses pièges. Le
chapitre~\ref{chap:donnees} présente les données, soit 45\,008 opérations
et la mesure de l'avantage qui leur est associée, puis chiffre la dépendance
qui menace les tests. On ne peut en effet juger proprement qu'après avoir
mesuré ce qui pourrait fausser le jugement.

Le chapitre~\ref{chap:juger} constitue le cœur du travail. On y juge les
stratégies par une escalade de prudence. Le test de référence est d'abord
appliqué sous ses hypothèses habituelles, puis les conditions qu'il suppose
sont mises à l'épreuve une à une, en neutralisant chacun des canaux de
dépendance mesurés au chapitre précédent. Le verdict n'est prononcé qu'au
terme de cette montée en exigence.

Le chapitre~\ref{chap:protocole} prend alors du recul. Il récapitule la
procédure sous la forme d'un outil réutilisable, en assume les limites et
trace les suites possibles. La conclusion referme l'enquête en répondant à la
question posée ici.

Le fil va ainsi du plus fragile, la donnée et sa dépendance, au plus solide,
le verdict puis l'outil qui permettra de juger les signaux à venir.
